\section{Introduction}
\label{sec:intro}

% ---------------------------------------------------------------------
% Paragraph 1: Motivation -- KT is the core of adaptive learning.
% ---------------------------------------------------------------------
Knowledge tracing (KT) is the task of modelling a learner's evolving
knowledge state from their sequence of interactions, and predicting
whether they will answer the next item correctly. It is the core of any
adaptive learning engine: the learner model drives the decision of what
to teach or assess next. This study is the research core of
\emph{AlgoVenture}, an adaptive, gamified platform for learning data
structures and algorithms; its adaptive engine estimates per-concept
mastery (UC6) and recommends the next exercise and difficulty (UC7), and
it adopts the model and policy that this study finds best. Rather than
collecting classroom data, we investigate KT for the
programming/algorithm-learning domain on large published datasets, which
keeps the study reproducible and the data-collection effort minimal.

% ---------------------------------------------------------------------
% Paragraph 2: The gap -- programming domain is under-studied.
% ---------------------------------------------------------------------
However, most KT research and public benchmarks target the mathematics
domain (e.g., ASSISTments), while the programming/algorithm-learning
domain is comparatively under-studied~\cite{piech2015dkt,liu2023simplekt}.
Programming submissions carry signals absent from a simple
correct/incorrect label---the knowledge components (KCs) a problem
exercises and the structure of the code a student wrote.

% ---------------------------------------------------------------------
% Paragraph 3: The deeper gap -- accuracy vs. utility.
% ---------------------------------------------------------------------
A second, deeper gap is methodological. KT research overwhelmingly
optimises next-step prediction accuracy (AUC), yet the link between
predictive accuracy and the quality of the downstream instructional
decision---which exercise to present next---is only weakly established
in general KT, and largely unexamined for programming. A model that
predicts responses well is not guaranteed to sequence exercises well.

% ---------------------------------------------------------------------
% Research question
% ---------------------------------------------------------------------
\paragraph{Research question}
For programming-concept learning, do KT models that exploit
programming-specific signals (knowledge-component tags and code
structure) (a)~predict learner mastery more accurately than
response-only models, and (b)~do those accuracy gains actually translate
into better next-exercise sequencing---and how well do the findings
generalise across public datasets and relative to the mathematics
domain?

% ---------------------------------------------------------------------
% Contributions
% ---------------------------------------------------------------------
\paragraph{Contributions}
This work makes the following contributions:
\begin{itemize}
  \item A reproducible programming-KT evaluation pipeline over
        $\geq 2$ public datasets with standard baselines
        (DKT, SAKT, AKT, simpleKT). \textbf{(RO1)}
  \item A programming-domain KT variant that integrates
        knowledge-component tags and code-structure features
        (AST / code embeddings) on top of a strong baseline. \textbf{(RO2)}
  \item An evaluation of predictive accuracy (AUC, RMSE) of the proposed
        variant against the baselines. \textbf{(RO3)}
  \item An evaluation of downstream next-exercise sequencing utility via
        offline / simulator-based evaluation, together with an analysis
        of how strongly predictive accuracy correlates with sequencing
        utility. \textbf{(RO4)}
  \item An assessment of cross-dataset and cross-domain (programming vs.
        mathematics) generalisation, released as an open-source
        reproducible benchmark. \textbf{(RO5)}
\end{itemize}

\paragraph{Scope}
The study uses published datasets only; no data is self-collected.
Beating every state-of-the-art KT model is a non-goal---the contribution
centres on the programming domain and the prediction-vs-utility question.

% ---------------------------------------------------------------------
% Purpose of the RBL research (academic / practical / community)
% ---------------------------------------------------------------------
\paragraph{Purpose of the research}
The RBL research within AlgoVenture serves three complementary purposes.
\begin{itemize}
  \item \textbf{Academic.} Verify whether integrating programming-specific
        signals (code structure and knowledge-component tags) into a
        knowledge-tracing model predicts student ability more accurately
        than response-only baselines, and---crucially---whether any such
        accuracy gain actually leads to better adaptive sequencing
        (choosing the next exercise).
  \item \textbf{Practical.} Rather than picking an off-the-shelf model
        arbitrarily, the pipeline acts as an empirical filter: the
        AlgoVenture platform \emph{consumes} the study's result by
        adopting the winning model (highest evaluation score) into its
        adaptive learning engine.
  \item \textbf{Scientific / community contribution.} Deliver an
        open-source benchmark pipeline that compares KT models
        transparently on public educational datasets (CodeWorkout,
        FalconCode). This lets the team produce a standards-compliant
        research report \emph{without} collecting any new classroom data.
\end{itemize}
